\chapter{Medicines (antidepressant)}
\index{Medicines (antidepressant)}

\section*{Definition and Overview}
Medicines (antidepressant), or antidepressants, are a diverse class of psychiatric pharmacological agents primarily prescribed for the management of moderate-to-severe depressive disorders, anxiety disorders, obsessive-compulsive disorder (OCD), post-traumatic stress disorder (PTSD), and chronic pain conditions. They act by modifying the concentration and activity of neurotransmitters---specifically serotonin, noradrenaline, and dopamine---within the synaptic cleft in the central nervous system. Modern treatment pathways emphasise that antidepressants are most effective when integrated into a comprehensive, individualised care plan that includes psychological therapies, lifestyle modifications, and close medical supervision.

\section*{Epidemiology}
The prescription rates of antidepressant medications have steadily risen in many developed nations. In many countries, antidepressants are among the most frequently prescribed classes of medication, with approximately 1 in 8 people taking an antidepressant daily. They are prescribed across a wide demographic, with significantly higher rates of use among females and older adults. Selective serotonin reuptake inhibitors (SSRIs) represent the most commonly prescribed subclass due to their favourable safety profile and tolerability compared to older agents.

\section*{Aetiology and Pathophysiology}
The therapeutic mechanisms of antidepressants target neurochemical pathways, receptor sensitivity, and neuroplasticity.
\begin{itemize}
    \item \textbf{Monoamine hypothesis:} The historic framework proposing that depressive symptoms are linked to a deficiency of monoaminergic neurotransmitters (serotonin, noradrenaline, and dopamine) in brain synapses.
    \item \textbf{Reuptake inhibition:} Transporter blockades (e.g.\ SERT and NET) prevent the reabsorption of serotonin and noradrenaline into the pre-synaptic neuron, increasing their extracellular availability in the synaptic cleft.
    \item \textbf{Receptor down-regulation:} Long-term administration leads to the desensitisation and down-regulation of pre-synaptic autoreceptors (like 5-HT1A), which subsequently increases post-synaptic neurotransmission.
    \item \textbf{Neurogenesis and brain-derived neurotrophic factor (BDNF):} Chronic antidepressant use stimulates BDNF expression, promoting neuroplasticity, dendritic branching, and neurogenesis in the hippocampus, which aligns with the typical 2-to-4-week delay in clinical efficacy.
    \item \textbf{Differentiation of subclasses:} Includes Selective Serotonin Reuptake Inhibitors (SSRIs, e.g.\ escitalopram), Serotonin-Noradrenaline Reuptake Inhibitors (SNRIs, e.g.\ venlafaxine), Tricyclic Antidepressants (TCAs, e.g.\ amitriptyline), and Monoamine Oxidase Inhibitors (MAOIs, e.g.\ phenelzine).
\end{itemize}

\section*{Clinical Features}
Antidepressant therapy is characterised by its delayed clinical response, potential initial side effects, and class-specific properties.
\begin{itemize}
    \item \textbf{Delayed therapeutic onset:} Clinical improvement in mood and anxiety typically takes 2 to 6 weeks to become noticeable, though somatic symptoms like sleep and energy may improve earlier.
    \item \textbf{Initial transient side effects:} Nausea, headache, mild agitation, and sleep disturbances are common during the first 1 to 2 weeks of starting therapy and usually resolve.
    \item \textbf{Sexual dysfunction:} Reduced libido, erectile dysfunction, and anorgasmia are common, long-term side effects associated particularly with SSRIs and SNRIs.
    \item \textbf{Weight changes:} Fluctuations in appetite and weight, with some agents (like mirtazapine) causing weight gain, while others (like fluoxetine) are weight-neutral or cause weight loss.
    \item \textbf{Discontinuation symptoms:} Dizziness, electric-shock sensations (``brain zaps''), anxiety, and flu-like symptoms if the medication is stopped abruptly rather than tapered.
\end{itemize}

\section*{Differential Diagnosis}
Before initiating antidepressant therapy, clinicians must rule out other underlying causes of depressive and somatic symptoms.
\begin{itemize}
    \item \textbf{Bipolar disorder:} Prescribing antidepressants without a mood stabiliser in bipolar patients can trigger manic or hypomanic episodes, necessitating careful screening.
    \item \textbf{Hypothyroidism:} Endocrine disorder presenting with fatigue, depressed mood, weight gain, and lethargy, diagnosed with a thyroid-stimulating hormone (TSH) test.
    \item \textbf{Vitamin D or B12 deficiency:} Nutritional deficiencies causing fatigue, cognitive changes, and mood disturbances, ruled out with routine blood chemistry.
    \item \textbf{Substance-induced mood disorder:} Depressive symptoms arising from alcohol abuse, illicit drug use, or side effects of other medications (e.g.\ corticosteroids or beta-blockers).
    \item \textbf{Obstructive sleep apnoea:} Often causes chronic fatigue, morning headaches, and low mood, which may be misdiagnosed as clinical depression.
\end{itemize}

\section*{When to Seek Medical Care}
Patients starting antidepressants or their families should maintain close contact with their prescribing doctor or mental health professional.

\textbf{Contact your local emergency services for:}
\begin{itemize}
    \item Sudden, severe worsening of suicidal ideation, self-harm behaviours, or planning suicide, especially in the first few weeks of therapy or after a dose change.
    \item Signs of serotonin syndrome, a life-threatening drug reaction characterised by high fever, severe agitation, muscle rigidity, tremor, hyperreflexia, and confusion.
    \item Severe, acute allergic reactions or signs of overdose (unresponsiveness, seizures, or cardiac arrhythmias).
\end{itemize}
Other reasons to see a doctor: See your doctor if side effects are persistent, or if you wish to change or stop taking your medication. Never stop antidepressants abruptly.

\section*{Diagnosis and Investigations}
Starting an antidepressant requires a comprehensive clinical assessment to select the most appropriate agent and establish baseline health.
\begin{itemize}
    \item \textbf{Mental state examination and history:} Assessing the severity of depression/anxiety, risk of self-harm, previous treatment responses, and family history.
    \item \textbf{Electrocardiogram (ECG):} Recommended before starting TCAs or specific SSRIs (like citalopram at higher doses) to screen for long QT syndrome or other cardiac conduction anomalies.
    \item \textbf{Electrolyte monitoring:} Measuring serum sodium levels, particularly in elderly patients, as SSRIs can cause the syndrome of inappropriate antidiuretic hormone secretion (SIADH) and hyponatraemia.
    \item \textbf{Renal and hepatic function tests:} Baseline blood tests to determine the starting dose and frequency, as many antidepressants are metabolised by the liver and cleared by the kidneys.
    \item \textbf{Blood pressure check:} Essential before and during treatment with SNRIs (like venlafaxine), which can cause dose-dependent increases in blood pressure due to noradrenergic activity.
\end{itemize}

\section*{Management}
Antidepressant management involves careful titration, compliance monitoring, side-effect management, and structured discontinuation.
\begin{itemize}
    \item \textbf{Structured initiation:} Starting at a low dose and gradually titrating upward to minimise side effects, while advising the patient of the expected delay in therapeutic benefits.
    \item \textbf{Combination therapy:} Combining pharmacotherapy with evidence-based psychotherapies (such as cognitive behavioural therapy - CBT) for optimal clinical outcomes.
    \item \textbf{Active monitoring:} Scheduling frequent appointments (e.g.\ weekly or fortnightly) during the initial phase of treatment to assess efficacy, safety, and compliance.
    \item \textbf{Treatment duration guidelines:} Continuing the medication for at least 6 to 12 months after achieving full remission of a first depressive episode to prevent relapse; recurrent episodes may require long-term maintenance.
    \item \textbf{Gradual tapering:} Discontinuing the medication by slowly reducing the dose over several weeks or months under medical supervision to avoid discontinuation syndrome.
\end{itemize}

\section*{Complications}
Although widely prescribed, antidepressants carry risks of adverse clinical events and interactions.
\begin{itemize}
    \item \textbf{Serotonin syndrome:} A potentially fatal toxic state caused by excess intrasynaptic serotonin, typically due to combining multiple serotonergic drugs (e.g.\ SSRIs with MAOIs, tramadol, or St John's wort).
    \item \textbf{Hyponatraemia:} Low sodium levels in the blood, presenting with confusion, muscle weakness, headache, and seizures, primarily in older patients taking SSRIs.
    \item \textbf{Increased suicidality:} A transient, paradoxical increase in suicidal thoughts and behaviours, particularly in children, adolescents, and young adults under 25, during the initial weeks of therapy.
    \item \textbf{QT prolongation and arrhythmias:} Specific agents (e.g.\ citalopram, tricyclics) can prolong the QT interval, predisposing patients to torsades de pointes and sudden cardiac arrest.
    \item \textbf{Bleeding risk:} SSRIs impair platelet serotonin uptake, increasing the risk of upper gastrointestinal bleeding, especially when co-administered with NSAIDs, aspirin, or anticoagulants.
\end{itemize}

\section*{Prognosis and Outcomes}
Approximately 50\% to 60\% of individuals with moderate-to-severe depression achieve a significant clinical response with their first antidepressant. For those who do not respond, switching to a different class, combining treatments, or adding an augmenting agent (such as an atypical antipsychotic or mood stabiliser) can achieve remission. When combined with psychotherapy, the long-term prognosis and prevention of future depressive episodes are significantly improved.

\section*{Prevention and Self-Care}
Self-care during antidepressant treatment involves understanding the medication and adopting healthy lifestyle practices.
\begin{itemize}
    \item \textbf{Take medication consistently:} Establish a daily routine to take the medication at the same time each day to maintain stable blood levels.
    \item \textbf{Avoid alcohol and recreational drugs:} These substances can worsen depression, increase sedative side effects, and decrease the efficacy of antidepressants.
    \item \textbf{Maintain open communication:} Keep a log of any side effects or changes in mood and share them with your prescribing doctor or psychologist.
    \item \textbf{Incorporate lifestyle measures:} Engage in regular aerobic exercise, practice good sleep hygiene, and eat a balanced diet, which are proven to support mental health recovery.
    \item \textbf{Build a support network:} Rely on family, friends, or support groups, and engage in therapy to develop healthy coping mechanisms.
\end{itemize}

\section*{Sources and Further Reading}
The following organisations provide further information and resources on antidepressant medications:
\begin{itemize}
    \item Beyond Blue. \textit{Support programs, fact sheets, and information on depression and medical treatments.} https://www.beyondblue.org.au/
    \item Black Dog Institute. \textit{Research-backed resources on depression, bipolar disorder, and antidepressant medications.} https://www.blackdoginstitute.org.au/
    \item NPS MedicineWise (nps.org.au / Healthdirect). \textit{Detailed medicine information, side effects, and safety tips for consumers.} https://www.healthdirect.gov.au/antidepressants
    \item Royal Australian and New Zealand College of Psychiatrists. \textit{Clinical practice guidelines for the treatment of mood disorders.} https://www.ranzcp.org/
    \item NHS. \textit{Comprehensive clinical overview, side effects, and safety information on antidepressants.} https://www.nhs.uk/mental-health/talking-therapies-medicine-treatments/medicines-and-psychiatry/antidepressants/
\end{itemize}
